% =====================================================================
% commands.tex — métadonnées, macros et environnements éditoriaux
% =====================================================================

% ---------- Métadonnées du rapport ------------------------------------
\newcommand{\universite}{Université de Béjaïa}
\newcommand{\departement}{Département Informatique}
\newcommand{\faculte}{[NOM DE LA FACULTÉ]}
\newcommand{\formation}{M1 ASR}
\newcommand{\module}{BDDA}
\newcommand{\anneeUniversitaire}{2026/2027}
\newcommand{\titreCourt}{Gestion d'une banque avec Oracle relationnel-objet}
\newcommand{\titreLong}{Gestion d'une banque (clients, comptes et opérations) avec Oracle (BD relationnel-objet)}
\newcommand{\etudiants}{[NOMS DES ÉTUDIANTS]}
\newcommand{\encadrant}{Dr. N. BOUADEM}
\newcommand{\datedepot}{02/12/2026}

% ---------- Macros de texte -------------------------------------------
\newcommand{\code}[1]{\texttt{#1}}
\newcommand{\ora}[1]{\texttt{#1}}
\newcommand{\terme}[1]{\textbf{#1}}
\newcommand{\important}[1]{\textbf{\color{navy}#1}}
\newcommand{\guil}[1]{\og{}#1\fg{}}
\newcommand{\capture}[1]{%
  \textcolor{grayText}{\footnotesize{}[Cf. capture n°~#1]}%
}
\newcommand{\annexe}[1]{Annexe~\ref{#1}}

% ---------- Questions / réponses de soutenance ------------------------
\newcommand{\qa}[4]{%
  \par\medskip
  \noindent\textbf{\color{navy}Q#1. #2}\par\nopagebreak
  \noindent\textit{Réponse courte.} #3\par\nopagebreak
  \noindent\textit{Si l'on insiste.} #4\par
}

% ---------- Encadrés éditoriaux ---------------------------------------
\newtcolorbox{clebox}[1][Points clés]{%
  enhanced,
  breakable,
  colback  = blue!4!white,
  colframe = navy,
  coltitle = white,
  fonttitle= \bfseries\sffamily,
  title    = {#1},
  boxrule  = 0.6pt,
  arc      = 2mm,
  left     = 3mm,
  right    = 3mm,
  top      = 2mm,
  bottom   = 2mm,
  before skip = 1.2em,
  after skip  = 1.2em
}

\newtcolorbox{attentionbox}[1][Point de vigilance]{%
  enhanced,
  breakable,
  colback  = orange!5!white,
  colframe = orange!70!black,
  coltitle = white,
  fonttitle= \bfseries\sffamily,
  title    = {\faExclamationTriangle~#1},
  boxrule  = 0.6pt,
  arc      = 2mm,
  left     = 3mm,
  right    = 3mm,
  top      = 2mm,
  bottom   = 2mm,
  before skip = 1.2em,
  after skip  = 1.2em
}

\newtcolorbox{remarquebox}[1][Remarque]{%
  enhanced,
  breakable,
  colback  = grayLight,
  colframe = grayText,
  coltitle = white,
  fonttitle= \bfseries\sffamily,
  title    = {\faInfoCircle~#1},
  boxrule  = 0.5pt,
  arc      = 2mm,
  left     = 3mm,
  right    = 3mm,
  top      = 2mm,
  bottom   = 2mm,
  before skip = 1.2em,
  after skip  = 1.2em
}

\newtcolorbox{objectifbox}[1][Objectif]{%
  enhanced,
  breakable,
  colback  = teal!5!white,
  colframe = teal,
  coltitle = white,
  fonttitle= \bfseries\sffamily,
  title    = {\faBullseye~#1},
  boxrule  = 0.6pt,
  arc      = 2mm,
  left     = 3mm,
  right    = 3mm,
  top      = 2mm,
  bottom   = 2mm,
  before skip = 1.2em,
  after skip  = 1.2em
}

% ---------- Raccourcis pour les listes --------------------------------
\setlist[itemize]{leftmargin=1.6em,itemsep=0.18em,topsep=0.25em}
\setlist[enumerate]{leftmargin=1.8em,itemsep=0.18em,topsep=0.25em}
